% Peekaboo-VL paper: 01_intro. Included from ../peekaboo_report.tex.
\section{Introduction}

\begin{figure*}[t]
\centering
\resizebox{\textwidth}{!}{%
\begin{tikzpicture}[node distance=5mm and 7mm]
% lane 1: video
\node[box, fill=io, minimum width=21mm] (clip) {UAV clip\\(any length, fps)};
\node[box, fill=io, right=of clip, minimum width=24mm] (cache) {Frame cache\\64 frames + fps\\+ frame indices};
\node[box, fill=vis, right=of cache, minimum width=30mm] (frames) {32 frames, 640$\times$352\\true timestamps\\\texttt{<1.9 s>}\,\texttt{<5.6 s>}\,\ldots\\(+ clock on timing Qs)};
\node[box, fill=vis, right=of frames, minimum width=24mm] (vit) {Vision tower (frozen)\\+ merger (LoRA)};
% lane 2: cue
\node[box, fill=cue, below=of frames, minimum width=30mm] (probe) {Motion probe (CPU)\\ORB $\to$ RANSAC similarity\\$\to$ pan / zoom / yaw};
\node[box, fill=cue, right=of probe, minimum width=24mm] (text) {Hint text, answer\\words: ``moving\\forward quickly\ldots''};
% lane 3: question
\node[box, fill=io, below=of probe, minimum width=30mm] (qa) {Question + options A--D};
% language model spanning the lanes
\node[box, fill=llm, right=9mm of text, minimum width=30mm, minimum height=34mm] (lm) {Language model\\(backbone)\\[2pt]reads frames,\\hint and question\\[2pt]LoRA adapters\\on its linear layers};
% outputs
\node[box, fill=llm, right=9mm of lm, minimum width=26mm, yshift=11mm] (head) {Next-token scores\\of A, B, C, D};
\node[box, fill=io, below=of head, minimum width=26mm] (dec) {PriDe / option-order TTA\\frame-time jitter\\(opt.) prior fusion};
\node[box, fill=io, below=of dec, minimum width=26mm] (out) {Answer letter};
\node[box, fill=cue, left=9mm of out, minimum width=24mm, yshift=-10mm] (blind) {Blind text scorer\\(no video)};
% edges
\draw[arr] (clip) -- (cache);
\draw[arr] (cache) -- (frames);
\draw[arr] (cache.south) |- (probe.west);
\draw[arr] (frames) -- (vit);
\draw[arr] (probe) -- (text);
\draw[arr] (vit.east) -- ++(2.5mm,0) |- ([yshift=11mm]lm.west);
\draw[arr] (text.east) -- (lm.west |- text.east);
\draw[arr] (qa.east) -- (lm.west |- qa.east);
\draw[arr] (qa.south) |- (blind.west);
\draw[arr] (lm.east |- head.west) -- (head.west);
\draw[arr] (head) -- (dec);
\draw[arr] (blind.east) -- ++(5mm,0) |- (dec.west);
\draw[arr] (dec) -- (out);
\end{tikzpicture}}
\caption{The Peekaboo-VL pipeline. Each clip is decoded once into a cache
that keeps its frame rate and which frames were taken, so the prompt can
carry true times. A CPU step measures how the camera moves between frames
and writes it as one line of text in the answers' own words. The backbone
language model, fine-tuned with LoRA, scores the four letters. The final step can remove letter bias
and, optionally, mix in a text-only prior.}
\label{fig:system}
\end{figure*}


Drone video is an unfriendly input for a video language model. The camera
never stops moving, often faster than anything on the ground, so ``the car
moves'' and ``the camera moves'' look alike. The cars are a few pixels
wide. And many questions ask \emph{when} something happens, which needs a
finer sense of time than the model usually gets.

The ReVA benchmark~\cite{yao2026reva} tests exactly this, with eleven
question tasks over clips from VisDrone, UAVDT, ERA and a new collection,
Hawk-UAV. The organizers' own model, ReMoSense, extends
Qwen2.5-VL-7B~\cite{qwen25vl2025} with learned camera-motion tokens and
object tokens and reaches 80.04\% on test. Their ablation says the
camera-motion tokens alone are worth 3.9 points overall and 9 points on
temporal questions. So camera motion matters. The question is how to get it
into the model.

Training new motion modules was not an option for us. So we asked a
simpler question: if the model cannot see the camera motion well, why not
just tell it? Drone footage makes that cheap. The ground is far away and
close to flat, so how one frame lines up with the next is described by a
shift, a zoom and a rotation. Classic computer vision measures those in a
few milliseconds on a CPU. We write the result as one line of text and let
fine-tuning decide how much to trust it (Fig.~\ref{fig:system}). The same
idea covers two other things the model cannot easily see: the true time of
each frame, which recent Qwen models read as text
timestamps~\cite{qwen3vl2025}, and the habits of whoever wrote the
answers, which a tiny text-only model picks up surprisingly well.

We entered the ReVA challenge with one A40 GPU and a few weeks.
This paper reports what that budget bought, including the experiments
that failed. Our contributions are:

\begin{itemize}
\item \textbf{Tell, don't learn.} We inject camera ego-motion as text
instead of adding learned motion tokens. A CPU probe measures the motion
between sampled frames in a few milliseconds, and we render it in the
benchmark's own answer vocabulary (``moving forward'', ``panning left''),
so the model can match it to the options. Fine-tuned with LoRA on
Qwen3.5-9B, the system reaches 83.6\% on the ReVA test leaderboard, 3.5
points above the organizers' purpose-built ReMoSense-7B, using no new
network parts and one GPU. We are explicit about where this margin comes
from: the 9B backbone alone already matches ReMoSense, and fine-tuning on
the seen videos adds the rest. After fine-tuning, the model demonstrably
reads the hint, but removing it costs only about 0.2 points overall.
\item \textbf{A benchmark audit.} The ReVA test set asks new questions
about videos that are already in training, and the official validation set
is largely a copy of training questions. We build a development split with
no training-answer leak that mirrors the test regime and tracks the test
score to about one point. A fifth of it are test questions whose answers
the validation set happens to contain, and we say so. We argue that ReVA
results should be reported alongside a trained blind baseline and per-task
numbers.
\item \textbf{Wording over measurement.} The same motion measurements,
written in image terms (``the image shifts up'') rather than the answers'
words, \emph{lowered} a small model's accuracy on camera-motion questions
by 12.5 points; the answer words raised it by 16. A text hint's wording is
first-order, not a detail.
\item \textbf{An empirical map, negative results included.} We measure
what each part is worth with pre-registered success rules: the motion hint
after fine-tuning, question packing, letter-prior debiasing and prior
fusion, an object-motion hint, a caption hint from the benchmark's own
question writer, a second training epoch, and frame density at test
time. Temporal grounding stays the hardest task, and its errors are
near-misses of about one second---the gap between the timestamps the model
sees.
\end{itemize}

The rest of the paper follows the system: data and split
(Sec.~\ref{sec:data}), what it is built on and how it compares to prior
work (Sec.~\ref{sec:related}, Sec.~\ref{sec:background}), the method
step by step (Sec.~\ref{sec:method}), the main results
(Sec.~\ref{sec:results}), and an ablation of every part
(Sec.~\ref{sec:ablation}).
